\section{Decay after clearing denominators}\label{sec:decay}

Since $\delta_n^{1/n}\to\ep$, clearing denominators costs an
exponential factor $\ep^{a+k}$. We bound the $n$th-root limsup
of the forms below its reciprocal.

\subsection{The exponential profile}

Every grid summand through $2rn$ vanishes. Write a remaining
point as $2rn+t$, where $t\ge1/M$. The products give
\begin{align*}
 P_n(2rn+t)
 &=\prod_{\theta\in\Theta}2^{Ln}
             \frac{\Gamma((t+\theta)/2+Ln)}{\Gamma((t+\theta)/2)},\\
 \prod_{h=0}^n(2rn+t+2h)
 &=2^{n+1}
      \frac{\Gamma(((2r+2)n+t)/2+1)}{\Gamma((2rn+t)/2)}.
\end{align*}
For $x=t/n$, Stirling's formula gives the profile
\begin{equation}\label{eq:Phi}
 \begin{split}
 \Phi(x)={}&m\log K
   +\frac{mD}{2}\bigl((x+2L)\log(x+2L)-x\log x\bigr)\\
 &-\frac a2\bigl((x+2r+2)\log(x+2r+2)
                    -(x+2r)\log(x+2r)\bigr)+a\log2.
 \end{split}
\end{equation}
Here $0\log0=0$. Since $(a-mDL)+mDL-a=0$, both the
$n\log n$ and linear Stirling terms cancel.

For each fixed $A>0$, uniformly on $1/M\le t\le An$,
\begin{equation}\label{eq:uniform-Stirling}
 \log R_n(2rn+t)=n\Phi(t/n)+O_A(\log n).
\end{equation}
Only the gamma factors in the denominator of $P_n$ can have
bounded arguments. Uniformly for $t\ge1/M$ and
$\theta$ in the fixed finite set $\Theta$,
\[
 \log\Gamma((t+\theta)/2)
      =\frac t2\log(t/2)-\frac t2+O(\log(2+t)).
\]
For large $t$ this is Stirling's formula; on the remaining
compact range it follows from continuity. The other gamma
arguments grow proportionally to $n$. This proves
\eqref{eq:uniform-Stirling} even when $t/n\to0$.

\subsection{Differentiation and the infinite tail}

Use Cauchy's estimate on the circle of radius
$\varepsilon=1/(2M)$ around $2rn+t$. The numerator factors
are positive at the real center. If a factor has value $c$
there, its absolute-value ratio on the circle is at most
$1+\varepsilon/c$. Using $1+x\le\exp(x)$ and multiplying
these bounds gives a cost $\exp(O(\log n))$, since the sum of
the reciprocals is $O(\log n)$. The inverse denominator ratios cost
$\exp(O(1))$, since the factors are at least $2rn+t$ and their
reciprocal sum is $O(1)$. Hence, on the bounded scaled range,
\[
 \sup_{|z|=\varepsilon}|R_n(2rn+t+z)|
                    \le n^C R_n(2rn+t)
\]
for a fixed $C$ and sufficiently large $n$.
Homogeneity gives $|W_n(2rn+t+z)|=O_A(n^E)$.
Cauchy's estimate
\[
 |V_n(2rn+t)|\le\varepsilon^{-k}
                  \sup_{|z|=\varepsilon}|T_n(2rn+t+z)|
\]
therefore shows that multiplication and differentiation cost
only a power of $n$. There are $O_A(n)$ terms in the bounded
range, so their absolute sum has $n$th-root limsup at most
$\exp(\max\Phi)$. Here $m,B,r,a$ remain fixed.

For $t\ge An$ with $A$ sufficiently large, every numerator
factor on the circle is at most $3t$, each denominator factor
is at least $t$, and $|W_n|\le C't^E$. Thus
\[
 |V_n(2rn+t)|\le C''\kappa_n3^{mDLn}
                         t^{-(a-mDL)n-a+E}.
\]
Put $q_n=(a-mDL)n+a-E>1$. The odd summation grid has spacing
$2/b$. Including the intervening grid points gives an upper bound
on a grid of spacing $1/b$. A first term and an integral then give
\[
 \sum_{\substack{t\text{ on the odd grid}\\t\ge An}}t^{-q_n}
 \le(An)^{-q_n}+b\frac{(An)^{1-q_n}}{q_n-1}.
\]
Stirling's formula gives the tail $n$th-root bound
\[
 2^aK^m3^{mDL}\ep^{-(a-mDL)}A^{-(a-mDL)}.
\]
The factor $n^{a-mDL}$ from the factorials cancels the factor
$n^{-(a-mDL)}$ from the lower tail endpoint $An$. The bound tends
to zero as $A$ grows because $a>mDL$.
Together with the bounded range this proves
\begin{equation}\label{eq:uncleared-decay}
 \limsup_{\substack{n\to\infty\\n\text{ odd}}}|F_{n,b}|^{1/n}
       \le G:=\exp\!\left(\max_{x\ge0}\Phi(x)\right).
\end{equation}

\subsection{The criterion implies cleared decay}

The function $f$ in \eqref{eq:f} satisfies
$2\Phi'(x)=\log f(x)$. After multiplication by a positive
denominator, $(\log f)'$ has numerator
\[
 2\bigl((a-mDL)x^2+2L(a-mDL)x-mDL(2r)(2r+2)\bigr).
\]
This quadratic has one positive zero. Thus $\log f$ first
decreases and then increases. It starts at $+\infty$ and
tends to $0$ from below, with leading term $-2(a-mDL)/x$.
It has one zero $x_0>0$, where $\Phi$ has its global maximum.
Substituting $\Phi'(x_0)=0$ into \eqref{eq:Phi} yields
\begin{equation}\label{eq:G}
 G=K^m(x_0+2L)^{mDL}2^a
 \left(\frac{(x_0+2r)^{2r}}{(x_0+2r+2)^{2r+2}}\right)^{a/2}.
\end{equation}
If $f(u)<1$, then $x_0<u$. The ratio in parentheses decreases
for $x\ge0$, since its logarithmic derivative is
$-2x/((x+2r)(x+2r+2))$. Using $e_*>\ep$ in \eqref{eq:G}
therefore gives
\[
 (\ep^{a+k}G)^2<U<1.
\]
The prime number theorem and \eqref{eq:uncleared-decay} imply
\begin{equation}\label{eq:cleared-decay}
 \limsup_{\substack{n\to\infty\\n\text{ odd}}}
        |\delta_n^{a+k}F_{n,b}|^{1/n}<1.
\end{equation}
Every fixed integer combination of the cleared forms also
tends exponentially to zero.
